\oldpage[I]{82}
\begin{proposition}[1.1.10]
\label{I.1.1.10-fr}
\begin{enumerate}
  \item[\rm(i)] Lorsque $f$ parcourt $A$, les ensembles $D(f)$ forment une
  base de la topologie de $X$.
  \item[\rm(ii)] Pour tout $f\in A$, $D(f)$ est quasi-compact. En particulier
  $X=D(1)$ est quasi-compact.
\end{enumerate}
\end{proposition}

(i) Soit $U$ un ensemble ouvert dans $X$~; par définition, on a
$U=X-V(E)$ où $E$ est une partie de $A$, et
$V(E)=\bigcap_{f\in E}V(f)$, d'où $U=\bigcup_{f\in E}D(f)$.

(ii) D'après (i), il suffit de prouver que si $(f_\lambda)_{\lambda\in L}$
est une famille d'éléments de $A$ telle que
$D(f)\subset\bigcup_{\lambda\in L}D(f_\lambda)$, il existe une partie finie
$J$ de $L$ telle que $D(f)\subset\bigcup_{\lambda\in J}D(f_\lambda)$. Soit
$\mathfrak{a}$ l'idéal de $A$ engendré par les $f_\lambda$~; on a par
hypothèse $V(f)\supset V(\mathfrak{a})$, donc
$\mathfrak{r}(f)\subset\mathfrak{r}(\mathfrak{a})$~; comme
$f\in\mathfrak{r}(f)$, il existe un entier $n\geq0$ tel que
$f^n\in\mathfrak{a}$. Mais alors $f^n$ appartient à l'idéal $\mathfrak{b}$
engendré par une sous-famille finie $(f_\lambda)_{\lambda\in J}$, et on a
$V(f)=V(f^n)\supset V(\mathfrak{b})=\bigcap_{\lambda\in J}V(f_\lambda)$,
c'est-à-dire $D(f)\subset\bigcup_{\lambda\in J}D(f_\lambda)$.

\begin{proposition}[1.1.11]
\label{I.1.1.11-fr}
Pour tout idéal $\mathfrak{a}$ de $A$,
$\operatorname{Spec}(A/\mathfrak{a})$ s'identifie canoniquement au sous-espace
fermé $V(\mathfrak{a})$ de $\operatorname{Spec}(A)$.
\end{proposition}

On sait en effet qu'il y a correspondance biunivoque canonique, respectant la
structure d'ordre de l'inclusion, entre idéaux (resp. idéaux premiers) de
$A/\mathfrak{a}$ et idéaux (resp. idéaux premiers) de $A$ contenant
$\mathfrak{a}$.

Rappelons que l'ensemble $\mathfrak{N}$ des éléments nilpotents de $A$ (ou
\emph{nilradical} de $A$) est un idéal égal à $\mathfrak{r}(0)$, intersection
de tous les idéaux premiers de $A$
(\hyperref[0.1.1.1-fr]{0, 1.1.1}).

\begin{corollary}[1.1.12]
\label{I.1.1.12-fr}
Les espaces topologiques $\operatorname{Spec}(A)$ et
$\operatorname{Spec}(A/\mathfrak{N})$ sont canoniquement isomorphes.
\end{corollary}

\begin{proposition}[1.1.13]
\label{I.1.1.13-fr}
Pour que $X=\operatorname{Spec}(A)$ soit irréductible
(\hyperref[0.2.1.1-fr]{0, 2.1.1}), il faut et il suffit que l'anneau
$A/\mathfrak{N}$ soit intègre (ou, ce qui revient au même, que l'idéal
$\mathfrak{N}$ soit premier).
\end{proposition}

En vertu du cor. (\hyperref[I.1.1.12-fr]{1.1.12}), on peut se borner au cas où
$\mathfrak{N}=0$. Si $X$ est réductible, il existe deux parties fermées
$Y_1$, $Y_2$ distinctes de $X$ et telles que $X=Y_1\cup Y_2$, d'où
$\mathfrak{j}(X)=\mathfrak{j}(Y_1)\cap\mathfrak{j}(Y_2)=0$, les idéaux
$\mathfrak{j}(Y_1)$ et $\mathfrak{j}(Y_2)$ étant distincts de $(0)$
(\hyperref[I.1.1.5-fr]{1.1.5})~; donc $A$ n'est pas intègre. Inversement, si
dans $A$ il y a deux éléments $f\neq0$, $g\neq0$ tels que $fg=0$, on a
$V(f)\neq X$, $V(g)\neq X$ (puisque l'intersection des idéaux premiers de
$A$ est $(0)$), et $X=V(fg)=V(f)\cup V(g)$.

\begin{corollary}[1.1.14]
\label{I.1.1.14-fr}
\begin{enumerate}
  \item[\rm(i)] Dans la correspondance biunivoque entre parties fermées de
  $X=\operatorname{Spec}(A)$ et idéaux de $A$ égaux à leurs racines, les
  parties fermées irréductibles de $X$ correspondent aux idéaux premiers de
  $A$. En particulier, les composantes irréductibles de $X$ correspondent aux
  idéaux premiers minimaux de $A$.
  \item[\rm(ii)] L'application $x\mapsto\overline{\{x\}}$ établit une
  correspondance biunivoque entre $X$ et l'ensemble des parties fermées
  irréductibles de $X$ (autrement dit toute partie fermée irréductible de $X$
  admet un point générique et un seul).
\end{enumerate}
\end{corollary}

(i) résulte aussitôt de (\hyperref[I.1.1.13-fr]{1.1.13}) et
(\hyperref[I.1.1.11-fr]{1.1.11})~; et pour démontrer (ii), on peut, en vertu de
(\hyperref[I.1.1.11-fr]{1.1.11}), se borner au cas où $X$ est irréductible~;
alors, d'après la prop. (\hyperref[I.1.1.13-fr]{1.1.13}), il existe dans $A$
un plus petit idéal premier $\mathfrak{N}$, qui correspond donc à un point
générique
